% =====================================================================
% Приложение В. Глоссарий терминов.
% ~30 терминов. Сгруппированы по тематическим блокам.
% =====================================================================

\section{Глоссарий}
\label{app:glossary}

В приложении приведены определения ключевых терминов работы.
Термины сгруппированы по тематическим блокам --- мульти-агентные
системы и языковые модели, коммуникационные топологии, фазы
и маршрутизаторы, человек в контуре управления, метрики и
инфраструктура, эксперимент и статистика. В пределах каждого
блока термины упорядочены логически от общего к частному.

\subsection*{Мульти-агентные системы и языковые модели}

\begin{description}[leftmargin=*, style=nextline]
\item[\textbf{Большая языковая модель} (от англ. Large Language
   Model, LLM)] Глубокая нейронная сеть с большим числом параметров,
   обученная на текстовом корпусе для предсказания следующего
   токена. В настоящей работе языковые модели используются
   в качестве рабочих агентов, симулятора человека, маршрутизаторов
   и судьи.
\item[\textbf{Мульти-агентная система} (от англ. Multi-Agent
   System, MAS)] Программная система, в которой несколько
   автономных агентов взаимодействуют между собой и с внешней
   средой для совместного решения задач.
\item[\textbf{Агент}] Программная сущность в мульти-агентной
   системе с фиксированной ролью, системной подсказкой, набором
   доступных инструментов и собственным внутренним состоянием.
\item[\textbf{Системная подсказка}] Постоянная инструкция, задающая
   агенту его роль, ограничения и стиль ответов. Не изменяется
   в пределах одного запуска.
\item[\textbf{Скрэтчпад} (от англ. scratchpad)] Внутренний
   локальный контекст агента, в котором хранится история
   сообщений и промежуточных рассуждений. Является приватным
   для агента в отличие от общего состояния графа. Организован
   по политике скользящего окна с автоматическим сжатием
   при переполнении.
\item[\textbf{Рабочий агент} (от англ. worker)] Агент, непосредственно
   участвующий в решении задачи, --- планировщик, исследователь,
   исполнитель, критик или дебатер. Противопоставляется
   управляющим компонентам --- координатору, маршрутизатору,
   судье оценки.
\item[\textbf{Судья оценки качества} (от англ. LLM-as-judge)]
   Внешняя языковая модель из другого семейства весов, вызываемая
   для оценки качества финального ответа системы. Использует
   несколько независимых вызовов с агрегацией по правилу
   самосогласованности.
\item[\textbf{Симулятор человека}] Языковая модель в роли участника
   мульти-агентной системы, имитирующая поведение человека в
   контуре управления. Управляется системной подсказкой,
   зависящей от заданной роли --- координатор, рецензент, судья,
   равноправный участник или наблюдатель.
\item[\textbf{Состояние графа} (от англ. graph state)] Общий
   контейнер, объединяющий контексты всех агентов, текущую
   фазу решения, активную топологию, журнал переходов и
   накопленные сигналы. Передается между узлами при исполнении
   и подвергается редукции при параллельном исполнении.
\item[\textbf{Бюджет прогона} (от англ. budget)] Верхний предел
   стоимости вызовов языковых моделей, выделенный на один
   прогон. Реализован в виде трехуровневого ограничителя ---
   на отдельный вызов, на полный прогон и на грид-эксперимент
   в целом.
\end{description}

\subsection*{Коммуникационные топологии}

\begin{description}[leftmargin=*, style=nextline]
\item[\textbf{Коммуникационная топология}] Направленный граф,
   описывающий допустимые маршруты обмена сообщениями между
   агентами мульти-агентной системы и порядок их активации.
\item[\textbf{Топология Star} (звезда)] Конфигурация, в которой
   центральный координатор последовательно вызывает специализированных
   работников и агрегирует их ответы.
\item[\textbf{Топология Chain} (цепочка)] Линейный конвейер из
   трех агентов с обратными связями для доработки.
\item[\textbf{Топология Mesh} (полносвязная)] Конфигурация с общей
   шиной сообщений, в которой все агенты обмениваются сообщениями
   со всеми по принципу циклической очереди.
\item[\textbf{Топология Debate} (дебатная)] Конфигурация, в которой
   два оппонента параллельно генерируют альтернативные решения,
   а судья выбирает победителя или синтезирует ответ.
\item[\textbf{Топология Hierarchical} (иерархическая)] Двухуровневая
   конфигурация из координатора верхнего уровня и двух подкоманд,
   работающих параллельно.
\item[\textbf{Адаптивная мета-топология}] Конфигурация, в которой
   активная базовая топология выбирается на каждой итерации
   маршрутизатором в зависимости от текущей фазы решения и
   наблюдаемых сигналов агентов. Авторская разработка настоящей
   работы.
\end{description}

\subsection*{Фазы и маршрутизаторы}

\begin{description}[leftmargin=*, style=nextline]
\item[\textbf{Фаза решения задачи}] Состояние конечного автомата
   процесса решения. Множество фаз содержит четыре значения ---
   планирование, исполнение, проверка и завершение.
\item[\textbf{Конечный автомат} (от англ. Finite State Machine,
   FSM)] Математическая модель вычислений с конечным множеством
   состояний и переходов между ними. В работе используется для
   моделирования последовательности фаз.
\item[\textbf{Сигнал агента}] Булев индикатор, выставляемый агентом
   в общем состоянии и наблюдаемый маршрутизатором. Примеры ---
   готовность к следующей фазе, признак тупикового состояния,
   счетчик отказов критика.
\item[\textbf{Маршрутизатор фаз} (\code{PhaseRouter})] Компонент,
   принимающий решение о переходе процесса в следующую фазу.
   Реализован в трех вариантах --- правиловом, на основе языковой
   модели и в виде оракула.
\item[\textbf{Маршрутизатор топологий} (\code{TopologyRouter})]
   Компонент, принимающий решение о смене активной базовой
   топологии внутри адаптивной мета-топологии.
\item[\textbf{Защитные правила переключения} (\code{SwitchGuards})]
   Набор ограничений, предотвращающих частое осцилляционное
   переключение топологий. Включает минимальное число итераций
   в одной топологии, интервал между сменами, лимит переключений
   за фазу и за весь запуск.
\item[\textbf{Ворота передачи состояния} (\code{TransitionGate})]
   Компонент, применяющий таблицу правил передачи общего
   состояния при смене активной топологии и регистрирующий факт
   перехода в журнале.
\end{description}

\subsection*{Человек в контуре управления}

\begin{description}[leftmargin=*, style=nextline]
\item[\textbf{Человек в контуре управления} (от англ. Human-in-the-
   Loop, HITL)] Парадигма, в которой автоматизированная система
   обращается к человеку за решением, оценкой или подтверждением
   в определенных точках процесса.
\item[\textbf{Шлюз взаимодействия с человеком} (\code{HumanGateway})]
   Программный протокол, через который агенты системы обращаются
   к человеку. Обеспечивает идемпотентность повторных вызовов по
   паре идентификаторов запуска и запроса.
\item[\textbf{Роль координатора}] Человек управляет планом и
   делегированием задач агентам. Активна в иерархических и
   полносвязных топологиях.
\item[\textbf{Роль рецензента}] Человек проверяет промежуточные
   результаты и инициирует их доработку. Естественно вписывается
   в Star и Chain.
\item[\textbf{Роль судьи}] Человек выносит финальный вердикт при
   наличии нескольких альтернативных решений. Применима в Debate.
\item[\textbf{Роль равноправного участника} (от англ. peer)]
   Человек участвует в обсуждении наравне с агентами, внося
   собственные сообщения в общую шину. Применима в Mesh.
\item[\textbf{Роль наблюдателя} (от англ. monitor)] Человек
   пассивно следит за процессом и вмешивается только в критических
   случаях --- превышение бюджета, зависание, циклы. Применима
   в любой топологии.
\item[\textbf{Маршрутизатор ролей} (\code{RoleRouter})] Компонент,
   динамически выбирающий активную роль человека в зависимости
   от фазы и наблюдаемого состояния.
\item[\textbf{Шкала NASA-TLX} (от англ. NASA Task Load Index)]
   Стандартный инструмент субъективной оценки когнитивной
   нагрузки по шести шкалам --- умственное, физическое и временное
   усилие, результативность, общее усилие, фрустрация.
\end{description}

\subsection*{Метрики}

\begin{description}[leftmargin=*, style=nextline]
\item[\textbf{Качество решения} ($q$)] Агрегированный показатель
   качества в диапазоне $[0, 1]$. Конкретная формула зависит от
   класса задач --- доля прошедших юнит-тестов, точное совпадение
   числа, среднее ROUGE-L с покрытием концептов.
\item[\textbf{Полная стоимость прогона} ($c$)] Сумма стоимостей
   всех вызовов языковых моделей за один прогон, выраженная в
   долларах США.
\item[\textbf{Стенное время прогона} ($\tau$)] Полное время
   выполнения одного прогона в секундах с момента старта до
   получения финального ответа.
\item[\textbf{Доля отказов} ($r_{\text{fail}}$)] Доля прогонов,
   завершившихся со статусом, отличным от \enquote{выполнено}.
   Включает превышение бюджета, превышение максимального числа
   итераций, ошибки вызова инструментов.
\item[\textbf{Число переключений топологии} ($N_{\text{switch}}$)]
   Фактическое число смен активной базовой топологии за один
   прогон адаптивной мета-топологии.
\item[\textbf{Доля заблокированных решений маршрутизатора}
   ($r_{\text{guard}}$)] Доля решений маршрутизатора топологий,
   заблокированных защитными правилами. Высокое значение
   свидетельствует о склонности маршрутизатора к частым
   переключениям.
\item[\textbf{Стоимостный вклад маршрутизатора} ($\gamma$)] Доля
   общего бюджета прогона, потраченная на собственные вызовы
   языковой модели маршрутизатором. Критическая метрика, отвечающая
   на вопрос, не превышает ли стоимость адаптации получаемый
   от нее выигрыш.
\item[\textbf{Доля регрессий качества} ($r_{\text{hurt}}$)] Доля
   задач, на которых адаптивная конфигурация дает качество хуже
   лучшей статической конфигурации. Ключевой индикатор
   безопасности механизма адаптации.
\item[\textbf{Разрыв до оракула} ($\Delta_{\text{LOO}}$)] Разность
   между качеством, обеспечиваемым оракулом по схеме исключения
   по одному, и качеством фактического маршрутизатора. Служит
   верхней границей достижимого качества при идеальной
   маршрутизации в рамках того же пространства решений.
\item[\textbf{Прокси-метрика когнитивной нагрузки}
   ($L_{\text{cog}}$)] Агрегированный показатель когнитивной
   нагрузки симулятора человека, объединяющий число обращений,
   объем предоставляемого контекста и среднюю задержку ответа.
   В экспериментах с реальными участниками заменяется на шкалу
   NASA-TLX.
\end{description}

\subsection*{Инфраструктура реализации}

\begin{description}[leftmargin=*, style=nextline]
\item[\textbf{LangGraph}] Библиотека на языке Python для построения
   и исполнения направленных графов состояния с поддержкой
   асинхронных вызовов, чекпойнтера и параллельного выполнения
   узлов.
\item[\textbf{Чекпойнтер}] Механизм сохранения промежуточного
   состояния графа в постоянное хранилище, позволяющий возобновлять
   прерванные запуски с последней успешной точки.
\item[\textbf{Apache Parquet}] Колоночный формат хранения больших
   объемов структурированных данных. В работе используется
   для журналов взаимодействия агентов и журналов вызовов
   языковых моделей.
\item[\textbf{Грид-эксперимент} (от англ. grid sweep)] Запуск
   серии прогонов по декартову произведению значений нескольких
   параметров с параллельным исполнением ячеек.
\item[\textbf{Прогон} (от англ. run)] Один полный цикл выполнения
   мульти-агентной системы от поступления задачи до получения
   ответа. Является элементарной единицей измерения в экспериментах.
\end{description}

\subsection*{Эксперимент и статистика}

\begin{description}[leftmargin=*, style=nextline]
\item[\textbf{Корпус задач}] Размеченное множество тестовых задач
   одного класса с заранее определенной процедурой автоматической
   оценки ответа. В работе используются четыре корпуса.
\item[\textbf{Воронка экспериментов}] Последовательность
   экспериментов, в которой каждый следующий сужает пространство
   рассматриваемых конфигураций по результатам предыдущего, что
   снижает суммарный объем вычислений без потери информативности.
\item[\textbf{Подтверждающий прогон} (от англ. confirmation run)]
   Дополнительный эксперимент на отличной от основной языковой
   модели, проводимый для проверки независимости выводов от
   выбора семейства весов.
\item[\textbf{Дисперсионный анализ} (от англ. analysis of
   variance, ANOVA)] Статистический метод проверки гипотезы о
   равенстве средних значений в нескольких группах путем
   сравнения межгрупповой и внутригрупповой дисперсии.
\item[\textbf{Поправка Бенджамини---Хохберга}] Процедура контроля
   доли ложных открытий при множественном статистическом
   сравнении.
\item[\textbf{Коэффициент эффекта Коэна} ($d$)] Стандартизированная
   мера величины различия между двумя средними. Используется
   для оценки практической значимости различий, а не только
   статистической.
\item[\textbf{Метод бутстрапа}] Непараметрический метод оценивания
   распределения статистики путем многократного ресэмплирования
   исходной выборки с возвращением. В работе применяется
   для построения доверительных интервалов средних значений
   метрик качества с уровнем доверия 95 процентов.
\item[\textbf{Доверительный интервал}] Интервал значений,
   накрывающий истинное значение параметра распределения с
   заданной вероятностью. В работе строится методом бутстрапа
   с уровнем доверия 95 процентов и тысячей повторений
   ресэмплирования.
\item[\textbf{Конфигурация системы}] Конкретный набор
   значений всех переменных эксперимента --- активная топология,
   роль человека, маршрутизаторы, языковая модель, корпус задач
   и значения сидов. Различают статические конфигурации (с
   фиксированной топологией) и адаптивные (с переключением
   топологии по фазам).
\item[\textbf{Метод исключения по одному} (от англ. leave-one-out,
   LOO)] Процедура построения верхнего потолка адаптации, в
   которой для каждой задачи оптимальная конфигурация вычисляется
   на остальных задачах того же класса.
\end{description}
